% !TeX root = ../../main.tex
% ------------------------------------------------------------------
%  4.2.5 Conexiones, puntos de falla y contingencia.
%
%  Puntos únicos de falla declarados conforme a RT-02.11.
% ------------------------------------------------------------------

\section{Conexiones, puntos de falla y contingencia}
\label{sec:conexiones}

Esta sección describe cómo se conectan los sitios, el terreno y la nube, identifica cada punto donde esa conexión o su infraestructura puede fallar, y declara cómo se resuelve cada falla o, cuando no se resuelve de forma automática, cuál es la contingencia. Su punto de partida es el caso: la conectividad actual de Puelche se corta con frecuencia y en algunos lugares no existe, por lo que la solución no puede suponer un enlace disponible y debe decir, para cada conexión, qué pasa cuando no lo está.

\subsection{Conexiones entre sitios, terreno y nube}

Los dos dominios fijos, la nube y el on-premise, se conectan por túneles VPN IPsec que terminan en el par de firewalls de cada sitio (D-01) y en el Transit Gateway de AWS en sa-east-1, con enrutamiento dinámico BGP. Cada sitio publica sus eventos directamente a la nube y opera sin conexión con los sistemas de otra instalación. Los caminos de acceso se distribuyen así:

\begin{itemize}
  \item Cada centro de distribución tiene tres caminos independientes: fibra D-03, LTE D-04 y Starlink D-06. Starlink permanece en espera caliente, con el terminal encendido, el túnel IPsec establecido y BGP de menor preferencia; solo toma el tráfico cuando fallan fibra y LTE y entonces prioriza DMS/WAL de Talca, broker y outbox, guías hacia ERP y SII, identidad y telemetría crítica.
  \item Cada cross-docking usa Starlink como camino principal y LTE de dos proveedores como respaldo.
\end{itemize}

El terminal satelital de los centros permanece encendido porque adquirir satélites y negociar el túnel durante una falla tardaría minutos; el plan de tarifa plana no agrega costo por mantenerlo encendido.

Los terminales EC55 y TC58e acceden por red celular a CloudFront, cuyas rutas \texttt{/v1} y \texttt{/sync/v1} llegan a la API REST pública.

Los flujos principales son las entradas HTTPS de usuarios externos por CloudFront, los túneles IPsec de los sitios hacia el Transit Gateway y la descarga local por Bluetooth de los termógrafos al terminal del conductor. La tabla siguiente detalla los medios, caminos alternativos y conmutación:

\begin{itemize}
  \item Termógrafo--terminal: los termógrafos descargan por Bluetooth al terminal del conductor.
  \item Nube--VM: AWS DMS accede únicamente a VM-02 de Talca por el túnel IPsec autenticado.
  \item Sitios--SQS: los eventos de cada sitio, incluidos los cross-docking, llegan a SQS FIFO sin depender de otro sitio.
\end{itemize}

Los 184 usuarios de administración, comercial y soporte de Talca y Concepción acceden a las consolas y a la administración de Keycloak solo por Verified Access. La consola Angular corre en dos tareas Fargate distribuidas entre dos zonas tras un ALB privado y reenvía las llamadas a la API REST privada por el endpoint de interfaz execute-api. CloudFront publica únicamente los endpoints OIDC necesarios para autenticación y renovación de sesión por la API REST pública sin autorizador hacia Keycloak tras el ALB privado; el verificador local de relevos permanece en cada sitio. Las cadenas usan el canal AS2 descrito en ADR-11 (Anexo 4-O), cuyo transporte mantiene dos tareas Fargate distribuidas entre dos zonas.

La Tabla~\ref{tab:conexiones} resume cada conexión con su medio, su camino alternativo y el tiempo de conmutación entre ambos.

\begin{tablalafrox}{Conexiones de la solución y su camino alternativo}{tab:conexiones}%
  {F{0.24} Y F{0.24} F{0.16}}%
  {\cab{Conexión} & \cab{Medio y protocolo} & \cab{Camino alternativo} & \cab{Conmutación}}
  CD Talca con la nube & Fibra D-03, túneles IPsec con BGP & LTE D-04 y, tras él, Starlink D-06 & Menos de 30 s \\
  CD Concepción con la nube & Fibra D-03, túneles IPsec con BGP & LTE D-04 y, tras él, Starlink D-06 & Menos de 30 s \\
  Cross-docking con la nube & Starlink D-06, túneles IPsec con BGP & LTE D-04 de dos proveedores & Menos de 30 s \\
  Terreno con la nube & Red celular, CloudFront, API REST pública & Almacén local del dispositivo & No aplica \\
  Termógrafos con el terminal del conductor & BLE con el terminal del conductor, en ruta y al regresar & Registro interno del termógrafo & No aplica \\
  Usuarios externos con los portales & Internet, HTTPS por CloudFront & Portal de Clientes instalado, que conserva el pedido hasta reconectar & No aplica \\
  Oficinas y trabajo remoto con las consolas & Internet, HTTPS por Verified Access & Otra conexión a Internet & No aplica \\
  Cadenas con el canal AS2 & Internet, TLS 1.3 al Network Load Balancer & Reenvío AS2 hasta el acuse MDN & No aplica \\
  Región primaria con la secundaria & Replicación nativa de Aurora, DynamoDB y S3 & Respaldo inmutable & Promoción autorizada \\
\end{tablalafrox}

{\footnotesize Fuente: elaboración propia.\par}

Los cinco sitios tienen caminos alternativos hacia la nube. Ninguna instalación depende de los sistemas de otra; el terreno opera sin señal y los termógrafos conservan su registro interno. Los portales requieren conexión, salvo el Portal de Clientes instalado, que conserva el catálogo descargado y el pedido armado hasta recuperar señal.

\subsection{Superficie de exposición}
\label{sub:superficie}

La Tabla~\ref{tab:superficie} delimita las únicas entradas desde redes externas. Cada entrada usa un subdominio propio del dominio del CLIENTE gestionado en Route 53.

\begin{tablalafrox}{Superficie de exposición}{tab:superficie}%
  {F{0.17} F{0.26} Y F{0.28}}%
  {\cab{Entrada} & \cab{Servicio y puerto} & \cab{Quién accede} & \cab{Control}}
  Portales y API & CloudFront, HTTPS 443 & Clientes, transportistas, proveedores y terminales de terreno & WAF, protección contra bots, Shield Advanced y autorizador \\
  Identidad & OIDC de Keycloak por CloudFront, HTTPS 443 & Personas usuarias & WAF y límites de tasa \\
  Consolas & Verified Access, HTTPS 443 & Oficina y trabajo remoto & Identidad, MFA y postura \\
  Canal AS2 & Network Load Balancer, TCP 443 con TLS 1.3 & Cadenas registradas & Lista de IP, certificados, firma y MDN \\
  Túneles de sitio & VPN del Transit Gateway, IPsec UDP 500 y 4500 & Firewalls D-01 de los cinco sitios & Autenticación IKE y cifrado IPsec \\
\end{tablalafrox}

{\footnotesize Fuente: elaboración propia.\par}

La tabla excluye toda otra entrada desde fuera de la red del CLIENTE. En recuperación ante desastres se activan las mismas entradas en us-east-1. MDM, EDR, SII, Transbank, GIS, notificaciones y telemetría de flota son dependencias de salida, detalladas en el Formulario T-11.

\subsection{Puntos de falla en los sitios y en los enlaces}

Cada conexión se apoya en equipos del sitio que también pueden fallar. La Tabla~\ref{tab:fallas-sitios} recorre esos puntos de falla, desde el enlace hasta la energía del recinto, e indica para cada uno el mecanismo que la resuelve y la contingencia si ese mecanismo no basta.

\begin{tablalafrox}{Puntos de falla en los sitios y en los enlaces}{tab:fallas-sitios}%
  {F{0.24} Y Y}%
  {\cab{Punto de falla} & \cab{Resolución} & \cab{Contingencia}}
  Fibra y LTE de un centro de distribución & BGP conmuta al túnel IPsec ya establecido por Starlink D-06, con prioridad para DMS/WAL, broker, guías, identidad y telemetría crítica. & La operación local conserva 24 h de autonomía. \\
  Caída del enlace de Talca para la oficina & Verified Access admite otra conexión a Internet con los mismos controles. & Plan de rutas y manifiesto aprobados en el WMS antes de las 22:00; tableros y costo de servir esperan el retorno. \\
  Starlink de un cross-docking & El LTE de dos proveedores toma el tráfico en menos de 30 s. & La ventana de 3 h opera 100\,\% local y sincroniza al reconectar. \\
  Enlace de Los Ángeles en la madrugada & Starlink queda como camino único, porque la red móvil falla a esa hora. & La ventana de madrugada opera íntegramente local. \\
  Firewall de Talca & La unidad pasiva asume las direcciones y los túneles en menos de 30 s. & Operación local 24 h sin enlace. \\
  Firewall de Concepción & La unidad pasiva asume el túnel y la conmutación de enlace en menos de 30 s. & La unidad dañada se reemplaza con la de reserva de Talca. \\
  Nodo del clúster de Talca & Las VM se reinician en los nodos restantes, con sus datos replicados por Ceph. & Ante la pérdida del clúster, se levanta el WMS de Talca en la nube sobre su copia en Aurora. \\
  Instancia de escritura VM-02 & Se reinicia en otro nodo del clúster. & D-05 restaura la base en hasta 4 h sin enlace. \\
  Disco de un nodo de Talca & Ceph marca el disco fuera y recompone sus réplicas en los otros nodos. & Durante la recomposición quedan dos réplicas. \\
  Miembro del stack de switches de Talca o de Concepción & El otro miembro mantiene el tráfico. & Reposición sin corte con la unidad de reserva de Talca. \\
  Servidor de borde de Concepción & RAID 10 y fuentes redundantes. & Reposición del equipo y reconstrucción de su base desde el estado central. \\
  Mini-PC de un cross-docking & Dos fuentes, dos SSD en RAID 1 y dos switches. & Reposición desde Talca y reconstrucción desde el estado central y SQS FIFO. \\
  Firewall de un cross-docking & La unidad pasiva del par asume el túnel en menos de 30 s. & La ventana de 3 h opera 100\,\% local y sincroniza al reconectar. \\
  Periféricos de andén & Sin redundancia por equipo. & La etiqueta se emite en otra de las impresoras del centro de distribución. \\
  Energía de la sala de Talca & UPS N+1 de 30 min y generador que toma carga en 8 a 15 s. & Estanque de 24 h y contrato de reabastecimiento. \\
  Corte de energía en la bodega & Las UPS de los gabinetes de piso sostienen switches PoE, AP e impresoras durante 30 min. & Registro de frío con generador en Talca y UPS de 30 min en Concepción. \\
  Climatización de la sala de Talca & Segunda unidad de precisión en N+1. & Alerta del monitoreo de los sensores ambientales. \\
\end{tablalafrox}

{\footnotesize Fuente: elaboración propia.\par}

En Concepción, el RAID 10 tolera la falla de un disco y las fuentes redundantes sostienen el servidor ante la pérdida de un circuito; si falla el equipo, su base se reconstruye desde el estado central alimentado por sus eventos. En Talca, el tablero respaldado por generador mantiene el registro de frío, y en Concepción la UPS del gabinete permite el apagado ordenado. La tabla muestra además que los firewalls van en par en los cinco sitios conforme a RT-08.03 (Bases Técnicas Transversales, cap. 8, p. 18). Concepción dispone de fuentes redundantes, y cada cross-docking tiene dos switches industriales; su mini-PC sigue siendo un equipo único por sitio, respaldado por alimentación redundante, discos en RAID 1, reposición desde Talca y reconstrucción desde la nube. La liberación sanitaria y el despacho con guía preemitida no dependen de la nube. En la ventana de Talca, VM-02 se reinicia en otro nodo y otra impresora emite la etiqueta si falla la del andén. Los mecanismos de alta disponibilidad se detallan en la sección~\ref{sub:3-alta-disponibilidad}.

\subsection{Puntos de falla en la nube, el terreno y las integraciones}

La Tabla~\ref{tab:fallas-nube} completa el recorrido con los puntos de falla que están fuera de los sitios: la nube, el terreno y los sistemas de terceros.

\begin{tablalafrox}{Puntos de falla en la nube, el terreno y las integraciones}{tab:fallas-nube}%
  {F{0.24} Y Y}%
  {\cab{Punto de falla} & \cab{Resolución} & \cab{Contingencia}}
  Zona de disponibilidad de sa-east-1 & Aurora conmuta en menos de 30 s, ElastiCache en menos de 60 s y Fargate reprograma las tareas. & La operación continúa en las zonas restantes. \\
  Región sa-east-1 completa & Tras autorización del CLIENTE se promueve Aurora, se restituyen identidad y API y se valida; entonces Route 53 dirige el tráfico a us-east-1. & RTO de 4 h y RPO de 15 min; bodega y terreno operan localmente durante la conmutación. \\
  Túnel VPN de Talca durante la replicación & BGP conmuta entre fibra D-03, LTE D-04 y Starlink D-06 para que DMS siga leyendo VM-02. & La autonomía local conserva la operación mientras se restituye cualquier camino. \\
  IdP maestro o enlace de identidad & La caché de solo lectura, de 24 h, conserva los datos de identidad recibidos, y el verificador local habilita los relevos de turno con el manifiesto firmado y el PIN personal. & La credencial de turno dura hasta 8 h en bodega y 14 h en terreno, sin revocación remota durante el corte, y queda disponible la cuenta de último recurso (ADR-06, Anexo 4-O). \\
  Señal móvil en ruta & La aplicación opera contra el almacén local del dispositivo. & Sincroniza al reconectar dentro de los 10 min de la Tabla~\ref{tab:t29}. \\
  ERP legado & El trabajador \texttt{erp-sync} en VM-04 retiene las solicitudes en la cola local o en SQS y reintenta con la misma clave idempotente. & Ningún camión sale sin documento tributario válido emitido por el ERP. \\
  Perfil de API saturado & El balanceador reparte entre tareas y el escalado agrega tareas cuando los procesos PHP-FPM ocupados superan el 70 \%. & Limitación de tasa en API Gateway con mensaje explícito al usuario. \\
  Mensaje de reconciliación inválido o que falla & El consumidor rechaza el sobre de edición desconocida sin aplicarlo; tras cinco intentos pasa a la cola de mensajes fallidos. & Alarma inmediata y reproceso manual trazable; el grupo afectado no bloquea a los demás. \\
  Excursión térmica crítica sin enlace & El gateway bloquea el despacho en menos de 5 s en el borde. & Buffer de 24 h y alerta al reconectar. \\
  SII, cadenas o notificaciones & Reintento con espera creciente y cola de mensajes fallidos. & Reproceso desde la cola al recuperarse el tercero. \\
  Carga del peak de septiembre & Fargate escala automáticamente por perfil y Aurora ya está dimensionada para 3$\times$ sin cambio de instancia. & Limitación de tasa en API Gateway. \\
\end{tablalafrox}

{\footnotesize Fuente: elaboración propia.\par}

Las fallas de nube o enlace no detienen la preparación ni la entrega, que se confirman localmente. La pérdida de la región primaria deja sin servicio a los portales y a la analítica hasta promover la réplica. El despacho conserva la guía preemitida al cerrar la carga nocturna; si la carga cambia entre las 05:30 y las 07:00, la guía queda inválida y el ERP de Talca emite otra por cualquiera de sus tres caminos, a los que Concepción y los cross-docking llegan por los suyos. Si no hay camino hasta el ERP, el camión sale solo con la carga amparada por su guía vigente y el ajuste pasa a la ruta siguiente. Si la guía ya se anuló, la salida queda bloqueada hasta disponer de un documento válido; ningún camión sale sin guía válida.
